\documentclass[a4paper,12pt]{article}

\usepackage[utf8]{inputenc} 
\usepackage{amsmath}
\usepackage[a4paper, margin=0.9in]{geometry} % Adjust the margin here
\usepackage{multirow}
    

\section*{Instructions}
\begin{enumerate}
    \item The time allocated for the exam is \underline{120 minutes}.
    \item During the exam, you \underline{are not} allowed to use any class materials or ask your neighbor for help.
    \item Calculators \underline{are not} allowed. You may only have writing supplies and a transparent water bottle without labels on your desk.
    \item Use only pens or non-erasable ink pens for the test; erasable pens are not allowed.
    \item Write your answers \underline{legibly}. Illegible answers will not be scored.
    \item The maximum score for the test is \underline{4 points}.
    \item Each question is worth \underline{1 point}, but the questions may not require the same amount of time. Therefore, it is advisable to read the entire assignment before starting to solve it. 
    \item When grading the test, only the solutions provided on the distributed sheets will be considered.
    \item Always provide your solutions in \underline{separate} sheets.
    \vfill
\end{enumerate}

\section*{Question 1}
\begin{itemize}
    \item [\textbf{(a)}] Describe the simple harmonic oscillation problem in one dimension with a simple figure. Write down and solve its equation of motion step by step to obtain $x(t)=Ae^{-i\omega_0 t}+Be^{i\omega_0 t}$, where $A$ and $B$ are some arbitrary constants, $\omega_0$ is the natural oscillation frequency and $t$ is time. -- \textit{0.5 points}

    \item [\textbf{(b)}] We know that the equation $x(t)=Ae^{-i\omega_0 t}+Be^{i\omega_0 t}$ can be written in $x(t)=C\cos{\left(\omega_o t\right)}$ form (where $C$ is an arbitrary real constant) that describes the position of the block with respect to time. What would be $C$ for the initial condition $x(0)=x_0$. What is the total energy of the oscillator? -- \textit{0.25 points}

    \item [\textbf{(c)}] Let us assume that the motion is damped with the resistance force $F_f=-b\dot{x}$. Given that the equation of motion is $x(t)=e^{-\beta t}x_0\cos{\left(\omega_1 t\right)}$, how long does it take for the damped oscillator to halve the amplitude of the oscillation? (Givens: $\ln2\approx 0.7$, $m=~1\text{kg}$, $b=1.4~\text{kg/s}$) -- \textit{0.25 points}
\end{itemize}

\section*{Solution:}
\section*{Question 2}
Consider two identical pendulums, each consisting of a mass $m$ attached to a string of length $L$ suspended from the same point. The pendulums are initially at rest. The pendulum $A$ is released from a height $h$ relative to $B$, while the pendulum $B$ is initially stationary.

\begin{itemize}
    \item [\textbf{(a)}] Determine the velocity of the pendulum $A$ just before it collides with the pendulum $B$. ~-- \textit{0.5 points}
    \item [\textbf{(b)}] After the collision, the two pendulums stick together and swing as one system. Calculate the maximum height reached by the combined system after the collision. -- \textit{0.5 points}
\end{itemize}

\section*{Solution:}
\section*{Question 3}
A car of mass $m$ is moving at a speed $v$ around a curve with a radius $R$. The coefficient of static friction between the tires and the road is $\mu_s$, and the gravitational acceleration is $g$.

\begin{itemize}
    \item [\textbf{(a)}] Determine the maximum speed $v_{\text{max}}$ at which the car can navigate the curve safely to avoid slippage. -- \textit{0.5 points}
    
    \item [\textbf{(b)}] If the road is banked at an angle $\theta$ with respect to the horizontal, how would the maximum speed of the car change compared to (a) for a safe turn? -- \textit{0.5 points}
\end{itemize}

Note: Assume that the car does not experience any other external forces, such as air resistance.

\section*{Solution:}
\section*{Question 4}
A spring with stiffness $k$ and zero equilibrium length is stretched by a force $F_0$ up to a point $x_{\text{max}}$ where $F_0$ is balanced with the spring force. At the end of the spring, a mass $m$ is attached. 
\begin{itemize}
    \item [\textbf{(a)}] When the mass is released, what is the work done by the spring force on the mass until it reaches half of its initial displacement? -- \textit{0.5 points}
    \item [\textbf{(b)}] Determine the impulse of the spring force during this specific phase of the motion. -- \textit{0.5 points}
\end{itemize} 

\section*{Solution:}


\section*{Question 1}
\begin{itemize}
    \item [\textbf{(a)}] Describe the simple harmonic oscillation problem in one dimension with a simple figure. Write down and solve its equation of motion step by step to obtain $x(t)=Ae^{-i\omega_0 t}+Be^{i\omega_0 t}$, where $A$ and $B$ are some arbitrary constants, $\omega_0$ is the natural oscillation frequency and $t$ is time. -- \textit{0.5 points}

    \item [\textbf{(b)}] We know that the equation $x(t)=Ae^{-i\omega_0 t}+Be^{i\omega_0 t}$ can be written in $x(t)=C\cos{\left(\omega_o t\right)}$ form (where $C$ is an arbitrary real constant) that describes the position of the block with respect to time. What would be $C$ for the initial condition $x(0)=x_0$. What is the total energy of the oscillator? -- \textit{0.25 points}

    \item [\textbf{(c)}] Let us assume that the motion is damped with the resistance force $F_f=-b\dot{x}$. Given that the equation of motion is $x(t)=e^{-\beta t}x_0\cos{\left(\omega_1 t\right)}$, how long does it take for the damped oscillator to halve the amplitude of the oscillation? (Givens: $\ln2\approx 0.7$, $m=~1\text{kg}$, $b=1.4~\text{kg/s}$) -- \textit{0.25 points}
\end{itemize}

\section*{Question 2}
Consider two identical pendulums, each consisting of a mass $m$ attached to a string of length $L$ suspended from the same point. The pendulums are initially at rest. The pendulum $A$ is released from a height $h$ relative to $B$, while the pendulum $B$ is initially stationary.

\begin{itemize}
    \item [\textbf{(a)}] Determine the velocity of the pendulum $A$ just before it collides with the pendulum $B$. ~-- \textit{0.5 points}
    \item [\textbf{(b)}] After the collision, the two pendulums stick together and swing as one system. Calculate the maximum height reached by the combined system after the collision. -- \textit{0.5 points}
\end{itemize}

\section*{Question 3}
A car of mass $m$ is moving at a speed $v$ around a curve with a radius $R$. The coefficient of static friction between the tires and the road is $\mu_s$, and the gravitational acceleration is $g$.

\begin{itemize}
    \item [\textbf{(a)}] Determine the maximum speed $v_{\text{max}}$ at which the car can navigate the curve safely to avoid slippage. -- \textit{0.5 points}
    
    \item [\textbf{(b)}] If the road is banked at an angle $\theta$ with respect to the horizontal, how would the maximum speed of the car change compared to (a) for a safe turn? -- \textit{0.5 points}
\end{itemize}

Note: Assume that the car does not experience any other external forces, such as air resistance.

\section*{Question 4}
A spring with stiffness $k$ and zero equilibrium length is stretched by a force $F_0$ up to a point $x_{\text{max}}$ where $F_0$ is balanced with the spring force. At the end of the spring, a mass $m$ is attached. 
\begin{itemize}
    \item [\textbf{(a)}] When the mass is released, what is the work done by the spring force on the mass until it reaches half of its initial displacement? -- \textit{0.5 points}
    \item [\textbf{(b)}] Determine the impulse of the spring force during this specific phase of the motion. -- \textit{0.5 points}
\end{itemize} 

\end{document}
